\section{The Loss of God}

Before the sentence names the fire, it says \emph{Depart}. Saint Alphonsus Liguori, whose \emph{Preparation for Death} sets the last things before the Christian one consideration at a time, found the whole of hell in that first word: ``The entire, then, of the Hell of the damned consists in that first word of the sentence of their condemnation,---`Depart from me, you cursed.'\,''\endnote{Alphonsus Liguori, \emph{Preparation for Death}, Consideration XXVI, third point (Dublin: James Duffy, 1843), p.~324.} The saints did not look away from the fire, but the Church teaches that the chief punishment of hell is the loss of God, and the Fathers and Doctors teach with her that this loss is the deeper wound.

\subsection{Separation from the light}

Saint Irenaeus of Lyons, writing in the second century against the Gnostics, begins his account of hell with communion with God, the good of which the lost are deprived:

\begin{quote}
And to as many as continue in their love towards God, does He grant communion with Him. But communion with God is life and light, and the enjoyment of all the benefits which He has in store. But on as many as, according to their own choice, depart from God, He inflicts that separation from Himself which they have chosen of their own accord. But separation from God is death, and separation from light is darkness; and separation from God consists in the loss of all the benefits which He has in store. Those, therefore, who cast away by apostasy these forementioned things, being in fact destitute of all good, do experience every kind of punishment. God, however, does not punish them immediately of Himself, but that punishment falls upon them because they are destitute of all that is good. Now, good things are eternal and without end with God, and therefore the loss of these is also eternal and never-ending. It is in this matter just as occurs in the case of a flood of light: those who have blinded themselves, or have been blinded by others, are for ever deprived of the enjoyment of light. It is not, [however], that the light has inflicted upon them the penalty of blindness, but it is that the blindness itself has brought calamity upon them.\endnote{Irenaeus, \emph{Against Heresies} V.27.2, trans.\ A. Roberts and W. H. Rambaut, \emph{The Ante-Nicene Fathers}, vol.~1.}
\end{quote}

Irenaeus says three things. The separation is chosen: God gives the soul what it has preferred. The suffering comes from the privation itself, as the blind man's darkness comes from his blindness and not from the light, which still floods the world. And the loss is everlasting because the goods of God, which the soul has refused, do not end.

\subsection{Ten thousand hells}

Saint John Chrysostom, preaching on Christ's warning that the tree which brings forth no good fruit is cast into the fire, drew from it the distinction the schools would later name the pain of loss and the pain of sense, and he told his people which of the two was the greater:

\begin{quote}
And though it appear indeed to be some single judgment, the being burnt up, yet if one examine carefully, these are two punishments. For he that is burnt is also cast of course out of God's kingdom; and this latter punishment is more grievous than the other. Now I know indeed that many tremble only at hell, but I affirm the loss of that glory to be a far greater punishment than hell.
\end{quote}

\noindent Then, having tried to make his hearers feel the loss by the image of a father cast out from the presence of a beloved child, he went on:

\begin{quote}
No doubt hell, and that punishment, is a thing not to be borne. Yet though one suppose ten thousand hells, he will utter nothing like what it will be to fail of that blessed glory, to be hated of Christ, to hear ``I know you not,'' to be accused for not feeding Him when we saw Him an hungered. Yea, better surely to endure a thousand thunderbolts, than to see that face of mildness turning away from us, and that eye of peace not enduring to look upon us.\endnote{Chrysostom, \emph{Homilies on Matthew} 23.9, \emph{Nicene and Post-Nicene Fathers}, 1st series, vol.~10.}
\end{quote}

The face Chrysostom speaks of is Christ's, and its turning away is what the word \emph{Depart} means. It is the face that sought the sinner through all his life, that was spat upon for him and looked on him from the cross, and that now, because he would not have it, will not look on him again.

Saint Augustine, who defended the eternity of punishment at length in the \emph{City of God}, says the same in his handbook of the faith. Even if the anger of God were the least that could be conceived, he writes, the loss itself would outweigh every pain:

\begin{quote}
to be lost out of the kingdom of God, to be an exile from the city of God, to be alienated from the life of God, to have no share in that great goodness which God has laid up for them that fear Him, and has wrought out for them that trust in Him, would be a punishment so great, that, supposing it to be eternal, no torments that we know of, continued through as many ages as man's imagination can conceive, could be compared with it.\endnote{Augustine, \emph{Enchiridion} 112, trans.\ J. F. Shaw, \emph{Nicene and Post-Nicene Fathers}, 1st series, vol.~3, New Advent, accessed 16 October 2026.}
\end{quote}

\subsection{An infinite loss}

Saint Thomas Aquinas gives the reason in the language of the schools. Every mortal sin, he teaches, has two sides: a turning away from God, the unchangeable good, and a disordered turning toward some changeable good. The punishment answers to each. ``Accordingly, in so far as sin consists in turning away from something, its corresponding punishment is the \emph{pain of loss,} which also is infinite, because it is the loss of the infinite good, i.e.\ God. But in so far as sin turns inordinately to something, its corresponding punishment is the \emph{pain of sense,} which is also finite.''\endnote{Thomas Aquinas, \emph{Summa theologiae} I-II, q.~87, a.~4, corpus and \latin{sed contra}, trans.\ Fathers of the English Dominican Province (Project Gutenberg eBook 17897).} The pain of sense is measured, and so it can be greater for one sinner and less for another; were the punishment infinite in every respect, Aquinas observes, ``the punishments of all mortal sins would be equal; because one infinite is not greater than another,'' and the councils teach that the punishments are unequal. The \emph{Catechism} likewise calls eternal separation from God ``the chief punishment of hell.''

Alphonsus puts the doctrine in plain words. ``But all these torments are nothing, compared with the pain of loss. Hell does not consist in the darkness, stench, shrieks and fire; the pain which constitutes Hell, is the pain of having lost God.'' A man who loses a jewel grieves according to its worth; the damned soul ``has lost God, who is an infinite good.''\endnote{Alphonsus, \emph{Preparation for Death}, Consideration XXVI, third point, p.~323.} And he imagines what we should hear if we could ask one of the lost why he weeps: ``I weep because I have lost God, and shall never see him more.''\endnote{Ibid., p.~325.} Saint Catherine of Siena heard the Eternal Father say the same in her \emph{Dialogue}: of the four principal torments of the damned, the first ``is, that they see themselves deprived of the vision of Me, which is such pain to them, that, were it possible, they would rather choose the fire, and the tortures and torments, and to see Me, than to be without the torments and not to see Me.''\endnote{Catherine of Siena, \emph{The Dialogue}, ``A Treatise of Discretion,'' the chapter ``Of the four principal torments of the damned,'' trans.\ Algar Thorold, new and abridged edition (London: Kegan Paul, 1907), Christian Classics Ethereal Library, accessed 16 October 2026.}

\subsection{Why we do not feel it now}

It is hard for us to fear this loss as the saints feared it, and Alphonsus says why: ``because they live in the midst of darkness, sinners who are content to live for months and years without God, do not understand this pain. However, they shall know at death the great good which they lose.''\endnote{Alphonsus, \emph{Preparation for Death}, Consideration XXVI, third point, p.~324.} In this life the soul is crowded with lesser goods, and a man can lose God for a season and hardly notice, because something else is always at hand to fill the place. At death those goods are taken from him. The soul, made for God, is left alone with its desire for him, and that desire becomes the measure of its anguish.

The pain of loss is felt in proportion to love, and so the saints, who loved God most, dreaded it most. Alphonsus records the saying of Ignatius of Loyola: ``Lord, I am willing to bear every pain, but not the pain of being deprived of you.''\endnote{Ibid., pp.~323--324.} Saint Francis de Sales sets the meditation on hell before the soul he guides toward devotion, and in the affections that close it he puts this loss to that soul as a question: ``whether you can bear to lose your God for ever?''\endnote{Francis de Sales, \emph{Introduction to the Devout Life} I.15, ``Seventh Meditation: On Hell,'' Affections and Resolutions (London: Rivingtons, 1876), p.~42.} Chrysostom answered his own warning with a prayer, and his prayer may be ours: ``But may this never be our case, O thou only-begotten Son of God, neither may we ever have any experience of this irremediable punishment.''\endnote{Chrysostom, \emph{Homilies on Matthew} 23.9.}
